\documentstyle[%


righttag%


,11pt%


,amssymb%


,russian%               remove in Eng.


,izvru%                 Set izven% in Eng.


]{amsart}





%       Math definitions





\newcommand{\la}{\langle}


\newcommand{\ra}{\rangle}


\newcommand{\tts}[1]{}





\theoremstyle{plain}


\theorembodyfont{\it}


\newtheorem{thmnonum}{\hskip\parindent Теорема}


\renewcommand{\thethmnonum}{}





%\hoffset=-15mm%                  For Eng.


 


\setcounter{page}{94}


\god{1998}


\nomer{11~(438)} %                        Remove in Eng.


%\nomer{Vol.\,42, No.\,11}%               Set in Eng.


%\str{\pageref{begin}--\pageref{end}}%  Set in Eng.


\udk{512.533.5}


\author{\it В.В.\,Бутаков, Л.М.\,Мартынов}


% NB:   ^^ remove in Eng.


\title{Полугруппы, все нециклические подполугруппы \\


которых являются идеалами}





\date{}


%\pagestyle{myheadings}%         Set in Eng.





\pagestyle{plain} %              Remove in Eng.


\begin{document}





%\thispagestyle{plain}%          Set in Eng.





\maketitle


\label{begin}





Изучение алгебр, подалгебры которых обладают тем или иным фиксированным


свойством, давно является одним из популярных направлений в алгебраических


исследованиях. В теории полугрупп указанный подход начал проявляться с


пятидесятых годов и разрабатывался многими авторами. Упомянем только


некоторые работы, инициировавшие наше исследование. Так, в статьях


\cite{1}--\cite{3} авторы независимо друг от друга описывают


полугруппы, в которых любая подполугруппа является идеалом


(здесь и в дальнейшем под идеалом понимается двусторонний идеал).


С другой стороны, в работе \cite{4} с точностью


до групп охарактеризованы полугруппы, все собственные подполугруппы


которых циклические. Условимся полугруппы с первым свойством называть


$Id$-{\it полугруппами}, а со вторым --- $C$-{\it полугруппами}.


Получив хорошее описание алгебр с некоторым свойством для всех подалгебр,


естественным является желание расширить этот класс, что зачастую


достигается требованием выполнимости зафиксированного свойства


лишь для некоторых подалгебр.





В данной работе изучаются полугруппы, все нециклические подполугруппы


которых являются идеалами; будем называть их $IdNC$-{\it полугруппами}.


Понятно, что любая $Id$-полугруппа и любая


$C$-полугруппа являются $IdNC$-полугруппами. В частности, любая


периодическая $C$-группа является таковой. Работа \cite{5} показывает,


что последние могут быть устроены весьма сложно. Используя описание


$Id$-полугрупп, получаем характеризацию $IdNC$-полугрупп.


При этом результат работы \cite{4} о $C$-полугруппах


не используется; мы получаем его в качестве следствия.





В работе используются терминология и обозначения, принятые в \cite{6}.


Подполугруппу $H$ полугруппы $S$ назовем $k$-{\it стянутой посредством


подмножества $M$ множества $S\setminus H$}, если $|M|\ge k$ и


$\la s_1, s_2, \dotsc, s_k\ra \supseteq H$ для любых различных элементов


$s_1, s_2, \dotsc, s_k$ из $M$. Если $Н$ является $k$-стянутой посредством


множества $S\setminus H$, то $H$ будем называть просто $k$-{\it стянутой}.


Обратим внимание на то, что термин ``стянутая полугруппа''


использовался в работе \cite{7} в другом смысле.





Определим некоторые типы полугрупп, которые будут использоваться при


описании $IdNC$-полугрупп.





{\it Tип} 1. Полугруппа $S$ есть $3$-нильпотентная, $S^2=\{c,0\}$,


$a^2=0$


для некоторого $a\in S\setminus S^2$ и $S^2$ является $2$-стянутой.





{\it Tип} 2. Полугруппа $S$ есть $3$-нильпотентная, $S^2=\{a^2, c, 0\}$,


подполугруппа $\{c,0\}$ является $2$-стянутой посредством множества


$S\setminus S^2$ и $\la x, y\ra \supseteq S^2$ для любых попарно различных


неразложимых элементов $x$ и $y$, из которых по крайней мере один


является делителем элемента $a^2$.





{\it Тип} 3. Полугруппа $S$ есть амальгама полугрупп $A$ и $B$ с общей


подполугруппой $C=A^3=B^2=\{c,0\}$, где $A$ есть раздувание циклической


нильпотентной полугруппы $\la  a\ra$ индекса $4$ над образующим $a$, причем


$c=a^3$,


$B$ есть нильпотентная $Id$-полугруппа ступени $3$ с $1$-стянутой


подполугруппой $C$, а операция умножения в $S$ совпадает на $A$ и $B$


с исходными операциями и для любых элементов $x\in A\setminus A^2$ и


$y\in B\setminus B^2$ имеет место $xy, yx \in C$ и $x^2y=yx^2=0$.





{\it Tип} 4. Полугруппа $S$ есть раздувание циклической нильпотентной


полугруппы индекса $4$ над образующим элементом.





{\it Tип} 5. Полугруппа $S$ есть раздувание циклической нильпотентной


полугруппы индекса $5$ над образующим элементом.





Напомним \cite{1}, что $Id$-полугруппы --- это в точности полугруппы,


в которых для любых элементов $x$ и $y$ cуществуют такие


натуральные числа $k$ и $r$, $2\le k,r\le 3$, что выполнено


условие $xy = x^k = y^r$. В частности, любая $Id$-полугруппа является


$3$-нильпотентной. Легко видеть, что ни одна из полугрупп указанных выше


типов 1--5 не является $Id$-полугруппой.





\begin{prop} Циклическая полугруппа $S = \la a\ra$ является


$IdNC$-полугруппой тогда и только тогда, когда она является конечной


индекса $r$ и периода $m$ и при этом выполнено одно из следующих


условий$:$


\begin{itemize}


\item[1)] $r\le 3$, $m$ любое$;$


\item[2)] $r=4$, $m$ нечетное$;$


\item[3)] $r=5$, $m$ не делится на $2$ и $3$.


\end{itemize}


\end{prop}





Доказательство предложения 1 проводится стандартными методами.





\begin{cor} Любая $IdNC$-полугруппа является периодической.


\end{cor}





\begin{prop} Нильполугруппа является $IdNC$-полугруппой тогда и


только тогда, когда она либо является $Id$-полугруппой, либо


является полугруппой одного из типов {\em 1--5}.


\end{prop}





При доказательстве этого утверждения сначала устанавливается, что


ступень нильпотентности любой нильпотентной $IdNC$-полугруппы


не превосходит 5, а затем с учетом того, что нильполугруппа, в которой


ступени всех ее нильпотентных подполугрупп ограничены в совокупности,


является нильпотентной (\cite {8}, теорема 3), делается вывод о


5-нильпотентности любой $IdNC$-нильполугруппы. Дальнейшее доказательство


использует стандартные комбинаторные методы.





Основным результатом работы является





\begin{thmnonum} Полугруппа $S$ является $IdNC$-полугруппой тогда и только


тогда, когда выполнено одно из следующих условий$:$





{\em I.} $S$ есть периодическая унипотентная полугруппа с максимальной


$C$-подгруппой $G$ такой, что любая собственная немаксимальная


подгруппа $H$ группы $G$ совпадает со своим идеализатором $I(H)$ в $S$,


а идеализатор $I(M)$ любой максимальной подгруппы $M$ из $G$ в $S$


есть циклическая $C$-полугруппа и $I(M)\cup G$ является идеалом в $S$,


а факторполугруппа Риса $S/G$ есть либо $Id$-полугруппа, либо


полугруппа одного из типов {\em 1--5;}





{\em II.} $S$ есть периодическая неунипотентная полугруппа и либо $S$ есть


правосингулярная $[$левосингулярная$]$ связка двух своих классов кручения


$K_e$, $K_f$, причем $K_eK_f=\{f\}$, $K_fK_e=\{e\}$


$[K_eK_f=\{e\}$, $K_fK_e=\{f\}]$,


либо $S$ есть $0$-прямое объединение полугрупп, одна из которых есть


$Id$-полугруппа, а каждая из остальных является циклической нильпотентной


полугруппой индекса $r\le 3$ с внешне присоединенным нулем.


\end{thmnonum}





Доказательство первого утверждения теоремы проводится стандартными


методами. Доказательство второго утверждения теоремы использует


теорему 2.6 из \cite{2}, в которой дано (в нашей терминологии) описание


$IdNC$-связок, и предложения 1 и 3 из \cite{9} о свойствах классов


кручения периодических полугрупп (можно использовать и более общие


результаты об эпигруппах из \S\,2 работы \cite{10}).





В нижеследующих двух случаях можно дать более детальное описание


соответствующих $IdNC$-полугрупп.





\begin{cor} Если $S$ --- унипотентная $IdNC$-полугруппа, $G$ --- ее


наибольшая подгруппа и $S/G$ является полугруппой типа $4$, то


$S$ является раздуванием циклической полугруппы индекса $4$ и


нечетного периода над образующим элементом.


\end{cor}





\begin{cor} Если $S$ --- унипотентная $INC$-полугруппа, $G$ --- ее


наибольшая подгруппа и $S/G$ является полугруппой типа $5$, то


$S$ является раздуванием циклической полугруппы индекса $5$ и периода,


не делящегося на $2$ и $3$, над образующим элементом.


\end{cor}





Следствием основного результата статьи является следующее





\begin{prop}[{\series{m}\shape{n}\selectfont{}\cite{4}}]


Полугруппа $S$ является $C$--полугруппой


тогда и только тогда, когда выполняется одно из следующих условий$:$


\begin{itemize}


\item[1)] $S$ --- двухэлементная полугруппа идемпотентов$;$


\item[2)] $S$ --- периодическая $C$-группа$;$


\item[3)] $S$ --- циклическая полугруппа либо индекса $2$ и любого периода,


либо индекса $3$ и нечетного периода$;$


\item[4)] $S$ --- раздувание конечной циклической группы над образующим


элементом некоторой ее максимальной подгруппы с помощью одного элемента$;$


\item[5)] $S$ --- раздувание конечной циклической группы над образующим


элементом с помощью двух элементов$;$


\item[6)] $S$ --- раздувание конечной циклической группы над двумя ее


различными образующими элементами с помощью одного элемента$;$


\item[7)] $S=\la a,b\ra$, где $a$, $b$ --- элементы индекса $3$ и нечетного


периода $m$, $a^k=b^k$ при $k\ge 2$ и $ab,ba\in \{a^2,a^{m+2}\}$.


\end{itemize}


\end{prop}





В заключение заметим, что $IdNC$-полугруппа не содержит в качестве


собственной подполугруппы квазициклическую $p$-группу


ни для какого простого числа $p$. С другой стороны, раздувание


периодической группы из \cite {5} над любым образующим


ее некоторой максимальной циклической подгруппы дает $IdNC$-полугруппу.





Авторы глубоко благодарны рецензенту за ряд полезных предложений


по улучшению первоначального текста статьи.





\begin{thebibliography}{99}





\bibitem{1} Шутов Э.Г. {\em Полугруппы с идеальными подполугруппами} //


Матем. сб. -- 1962. -- Т.\,57. -- \No\,2. -- С.\,179--186.


\bibitem{2} Miller D.W. {\em Hamiltonian semigroups} // Portugal Math. --


1962. -- V.\,21. -- \No\,3--4. -- P.\,171--187.


\bibitem{3} Kimura N., Tamura T., Merkel R.B. {\em Semigroups in which


all subsemigroups are left ideals} // Canad. J. Math. -- 1965.


-- V.\,17. -- \No\,1. -- P.\,52--62.


\bibitem{4} Баранский В.А., Нагребецкий В.Т. {\em Полугруппы, все


истинные подполугруппы которых циклические} // Тезисы докл. ХХVI конференция


матем. кафедр педиститутов Урала. -- Киров, 1968.


\bibitem{5} Ольшанский А.Ю. {\em Бесконечная группа с подгруппами простых


порядков} // Изв. АН СССР. -- 1979. -- Т. 44. --\ No\,2. -- C.\,309--321.


\bibitem{6} Шеврин Л.Н. {\em Полугруппы.} Гл.\,IV в кн. ``Общая алгебра''


/ Под ред. Л.А.\,Скорнякова. T.\,2. -- М.: Наука, 1991. -- С.\,11--191.


\bibitem{7} Просвиров А.С., Шеврин Л.Н. {\em Идеализаторы подполугрупп


и строение полугрупп } // Матем. зап. Уральск. ун-та. -- 1974.


-- Т.\,9. -- \No\,1. -- С.\,66--83.


\bibitem{8} Шеврин Л.Н. {\em О полугруппах, все подполугруппы которых


нильпотентны} // Сиб. матем. журн. -- 1961. -- Т.\,II. -- \No 6. --


C.\,936--942.


\bibitem{9} Просвиров А.С. {\em О периодических полугруппах} // Матем.


зап. Уральск. ун-та. -- 1971. -- Т.\,8. -- \No\,1. -- С.\,77--94.


\bibitem{10} Шеврин Л.Н. {\em К теории эпигрупп. } I // Матем. сб. -- 1994.


-- Т.\,185. -- \No\,8. -- C.\,129--160.


\end{thebibliography}





\vskip 0.5cm





\post{\it Омский государственный}{\it Поступила}


\post{\it педагогический университет}{07.07.1997}





\label{end}


\end{document}


